\documentclass[10pt]{article}

\usepackage[T1]{fontenc}
\usepackage{lmodern}
\usepackage{microtype}
\usepackage[margin=0.78in]{geometry}
\usepackage[numbers,sort&compress]{natbib}
\usepackage{multicol}
\usepackage{graphicx}
\usepackage[section]{placeins}
\usepackage[dvipsnames]{xcolor}
\usepackage{tikz}
\usepackage[
  colorlinks=true,
  linkcolor=MidnightBlue,
  citecolor=MidnightBlue,
  urlcolor=MidnightBlue,
  pdftitle={Molecular Design Software Is Now Free. Molecular Prediction Is Not.},
  pdfauthor={Rafal P. Wiewiora and Woody Sherman}
]{hyperref}

\setlength{\parindent}{0pt}
\setlength{\parskip}{0.48em}
\setlength{\bibsep}{0.5pt}
\setlength{\columnsep}{1.3em}
\renewcommand{\refname}{References}
\renewcommand{\bibfont}{\footnotesize}
\newcommand{\molarium}{\textsc{Molarium}}

\title{\vspace{-2.5em}\textbf{Molecular Design Software Is Now Free.\\Molecular Prediction Is Not.}}
\author{Rafal P. Wiewiora \and Woody Sherman}
\date{Opinion draft --- August 2026}

\begin{document}

\maketitle
\vspace{-1.4em}

\begin{abstract}
Computational molecular design has traditionally bundled scientific models, numerical implementations, interfaces, workflow plumbing, and expert practice into a single software product.
Coding agents are beginning to separate these layers.
They can increasingly construct interfaces, integrations, analyses, and bounded scientific tools on demand, sharply reducing the time and marginal cost required to reproduce implementation.
\molarium{} is our case study: a local-first workbench that runs substantial calculations in the browser.
It unifies molecular viewing, editing, preparation, conformer workflows, molecular mechanics, neural potentials, and protein-structure prediction.
Reference comparisons, model cards, provenance, a local security boundary, and explicit exclusions make its scope inspectable.
Accessibility is still not predictive validity.
We argue that defensible value is moving away from fixed feature inventories and toward validated scientific models, proprietary discovery intelligence, and reusable skills that encode how an organization judges evidence.
Molarium makes the tension concrete: the interface may become fluid; the scientific contract may not.
\end{abstract}

\section*{The interface is no longer the moat}

Historically, a request for a new molecular workflow entered a queue.
Scientists translated the question into requirements; engineers connected formats, programs, databases, and visualizations; product teams anticipated edge cases; and organizations maintained the result.
Even when an algorithm was public, turning it into dependable software required scarce time and specialized engineering.

Coding agents change the latency between a scientific request and an executable artifact.
Natural-language descriptions, examples, tests, and reference outputs can now act as a living specification.
An agent can inspect a repository, write and run code, compare outputs, diagnose failures, and revise the implementation.
Welsh predicted that conventional programming would recede behind learned systems~\cite{welsh2023end}.
The transition is no longer confined to programs encoded in model weights: the ordinary software surrounding scientific models is itself becoming generative.

Calling molecular design software ``free'' is therefore an economic claim, not a literal one.
Compute, support, security, maintenance, expert review, and validation still cost money.
What is collapsing is the scarcity of implementation: the lead time and marginal effort needed to reproduce a wrapper, interface, report, integration, or project-specific workflow.
A capability becomes commoditized before it becomes perfect.
It is enough that more people can rebuild and adapt it at lower cost.

This shift weakens a commercial moat based primarily on arranging public methods behind a polished interface.
Usability still matters, but interfaces become easier to imitate and personalize.
The durable application increasingly serves as a fluid presentation layer over scarcer assets: trusted models, curated data, prospective evidence, secure operations, and organization-specific judgment.

\section*{Molarium is a workbench, not a demo}

\molarium{} is an MIT-licensed, local-first browser workbench for molecular viewing, editing, preparation, simulation, and analysis~\cite{molarium2026}.
WebAssembly and WebGPU run most supported calculations on the user's device~\cite{haas2017webassembly,w3c2025webgpu}.
One explicit molecular state connects functions that otherwise span desktop applications, command-line programs, notebooks, and servers.

Molarium accepts SMILES, Protein Data Bank (PDB), XYZ, and prepared fixtures.
It displays proteins, ligands, waters, ions, hydrogen bonds, and aromatic contacts while preserving explicit topology.
It batches chemistry edits, validates the result, and relaxes the affected neighborhood.
Intermediate tautomers and bond edits never masquerade as finished states.
Undo, provenance, and structure export keep every intervention inspectable and reversible.

Protein preparation resolves models and alternate locations, then applies residue templates and connection records.
It rebuilds supported side-chain atoms, adds standard hydrogens, orients rotatable polar hydrogens, and screens crystallographic waters.
An audit records each change.
Preparation stops on unsupported chemistry, missing backbone atoms, and severe generated-atom clashes.
It does not build missing loops, metal coordination, covalent ligands, membranes, or periodic solvent boxes~\cite{molarium2026}.
The workbench acts only where its rules are defined.

Three task-specific views share the same state (Figure~\ref{fig:interface-workflow}).
View shows the protein, ligand, waters, contacts, and selected residue.
Build stages coupled bond-order, charge, and hydrogen edits and can optimize the ligand or local pocket without moving the rest of the protein.
Inspect shows preparation findings, residue context, and aligned trajectories; exports retain raw coordinates.
The controls change; molecular state and provenance persist.

\begin{figure}[htbp]
  \centering
  \includegraphics[width=\linewidth]{figures/molarium-interface-workflow.png}
  \caption{Three views of the same 7KPA complex: (A) View shows the protein, ligand, waters, and contacts; (B) Build stages chemistry and pocket-local edits; and (C) Inspect links a selected residue to preparation findings. Controls change; molecular context persists.}
  \label{fig:interface-workflow}
\end{figure}

The browser owns the molecular state, dispatches named engines, and returns evidence-bearing artifacts (Figure~\ref{fig:molarium-architecture}).
A new interface need not redefine the state, engine, or numerical scope.

\begin{figure}[htbp]
  \centering
  \begin{tikzpicture}[
    component/.style={
      draw=MidnightBlue,
      rounded corners=2pt,
      minimum height=1.05cm,
      text width=0.255\linewidth,
      align=center,
      inner xsep=5pt,
      font=\small
    },
    evidence/.style={
      draw=ForestGreen!70!black,
      rounded corners=2pt,
      minimum height=0.72cm,
      text width=0.80\linewidth,
      align=center,
      inner xsep=5pt,
      font=\small
    }
  ]
    \draw[dashed, draw=Gray] (-7.75,-1.22) rectangle (7.75,1.72);
    \node[anchor=west, font=\footnotesize, text=Gray] at (-7.58,1.48) {Browser-local scientific session};
    \node[component, fill=Gray!12] (state) at (-5.15,0.66) {\textbf{Scientific state}\\structures, explicit topology, protonation, user edits};
    \node[component, fill=SkyBlue!12] (workflow) at (0,0.66) {\textbf{Workflows}\\build, prepare, search, minimize, simulate, analyze};
    \node[component, fill=Cerulean!14] (engines) at (5.15,0.66) {\textbf{Versioned engines}\\RDKit, OpenMM/Sage, WebGPU, ANI-2x, OpenFold};
    \node[evidence, fill=ForestGreen!13] (evidence) at (0,-0.72) {\textbf{Evidence layer} --- model and asset hashes, reference comparisons, scope checks, preparation reports, structures, trajectories, and provenance};
    \draw[->, thick, MidnightBlue] (state.east) -- (workflow.west);
    \draw[->, thick, MidnightBlue] (workflow.east) -- (engines.west);
    \draw[->, thick, MidnightBlue] (-5.15,0.08) -- (-5.15,-0.34);
    \draw[->, thick, MidnightBlue] (0,0.08) -- (0,-0.34);
    \draw[->, thick, MidnightBlue] (5.15,0.08) -- (5.15,-0.34);
  \end{tikzpicture}
  \caption{Molarium's local-first architecture. Named engines act on explicit molecular state; evidence and provenance remain attached.}
  \label{fig:molarium-architecture}
\end{figure}

\section*{One workbench, several scientific contracts}

Each calculation route keeps its scientific identity.
RDKit handles chemical perception, ETKDGv3 embedding, and short MMFF94 or universal-force-field cleanup.
Browser-native small molecules combine OpenFF Sage valence and van der Waals parameters~\cite{boothroyd2023sage} with labeled Gasteiger charges.
Molarium can execute that numeric system through OpenMM's double-precision Reference platform compiled to WebAssembly~\cite{eastman2024openmm}, an experimental direct Sage WebGPU implementation, or a STORMM-style WebGPU engine for homogeneous replicas.
Comparisons reuse terms and coordinates; labels expose the engine and numerical scope.

The Conformer Arena exposes the distinction.
A worker pool creates shared RDKit seeds and briefly cleans them with MMFF94 or the universal force field.
Molarium compares the untouched seeds, a Sage/OBC2 STORMM refinement, and a separate ANI-2x refinement.
Every candidate receives a common Sage/OBC2 score; compatible candidates also receive a vacuum ANI-2x score~\cite{gao2020torchani}.
Users can compare landscapes and rank orderings, but Molarium does not subtract absolute energies from different Hamiltonians.
Symmetry-aware alignment changes display coordinates only; calculations and exports retain raw coordinates.

Protein support is narrower.
Browser engines execute preparameterized OpenFF Rosemary fixtures but do not parameterize arbitrary proteins.
A trained, template-free OpenFold export predicts one short protein chain after a configured service supplies a multiple-sequence alignment~\cite{ahdritz2024openfold}.
Its model card records the checkpoint, transformation, hashes, residue limits, recycle settings, and absence of templates or Amber relaxation.
These are separate scientific contracts.

Connected mode contacts external services only for explicit structure lookup or alignment search; calculation coordinates remain in the browser.
Local Lab applies a restrictive content-security policy, disables external structure and sequence controls, verifies reviewed executable-file hashes, and probes attempted egress with a synthetic canary~\cite{molarium2026}.
Users must still verify the checkout and control the machine.
The privacy boundary is testable, not absolute.

\section*{Molarium makes validation a product feature}

Molarium separates execution from evidence.
Its direct Sage and STORMM-style engines compare energies and forces with independent references on identical parameter arrays and coordinates.
Tests cover finite-difference forces, trajectory endpoints, constraint residuals, deterministic replay, replica isolation, model hashes, supported elements, buffer limits, and non-finite failures.
Native OpenMM, native TorchANI, hand-computable fixtures, and JavaScript double-precision calculations provide independent oracles~\cite{molarium2026}.
Unknown terms and unsupported domains fail closed.

Correct implementation is not valid prediction.
Evidence rises from execution, to agreement with declared equations, to retrospective benchmarks, prospective transfer, and improved experimental decisions.
A lower rung does not establish the next.
Time-split validation tests prospective performance more faithfully than random partitions~\cite{sheridan2013timesplit}; data quality, bias, and changing experimental context still constrain apparent progress~\cite{bender2021realistic}.

Figure~\ref{fig:evidence-ladder} places Molarium on the lower rungs.
Cross-backend comparisons, model cards, and scope checks support execution and implementation agreement, not prospective validity.
Higher rungs require temporally realistic benchmarks, predictions made before experiments, experimental feedback, and better decisions.

\begin{figure}[htbp]
  \centering
  \begin{tikzpicture}[
    stage/.style={
      draw=MidnightBlue,
      rounded corners=2pt,
      minimum height=0.48cm,
      text width=0.78\linewidth,
      align=left,
      inner xsep=6pt,
      font=\small
    }
  ]
    \node[stage, fill=ForestGreen!16] (s5) at (0,3.0) {\textbf{5. Improves decisions and outcomes} --- prospective evidence connects prediction to consequential experimental choices.};
    \node[stage, fill=RoyalBlue!14] (s4) at (0,2.25) {\textbf{4. Transfers prospectively} --- the model remains useful for new chemistry, targets, and experimental contexts.};
    \node[stage, fill=Cerulean!14] (s3) at (0,1.5) {\textbf{3. Performs retrospectively} --- a versioned benchmark measures accuracy, calibration, and domain coverage.};
    \node[stage, fill=SkyBlue!12] (s2) at (0,0.75) {\textbf{2. Matches its declaration} --- equations, parameters, and reference implementations agree on specified cases.};
    \node[stage, fill=Gray!12] (s1) at (0,0) {\textbf{1. Executes} --- the workflow runs, records provenance, and returns a technically valid artifact.};
    \draw[->, thick, MidnightBlue]
      (-7.1,0) -- (-7.1,3.0)
      node[midway, left=1mm, align=center, font=\footnotesize, text=MidnightBlue] {More scientific\\evidence};
  \end{tikzpicture}
  \caption{An evidence ladder for agentic molecular design. Automation can accelerate every rung, but evidence at a lower rung does not establish the next.}
  \label{fig:evidence-ladder}
\end{figure}

A preliminary conformer screen shows the distinction.
Across eight molecules, STORMM found the lower minimum for four, with a median molecule-level gain of only 0.128 kcal/mol relative to the RDKit seed pool.
It improved 41 of 224 paired starting conformers and added 75 clusters absent from the untouched seed lane; 44 clusters were within 3 kcal/mol in the STORMM lane but not in the untouched lane~\cite{molarium2026}.
Same-coordinate STORMM and OpenMM energies agreed within tolerance, so implementation mismatch did not explain the mixed result.
The evidence supports faster batch exploration and greater conformational diversity, not generally lower-energy conformers.

The limits are explicit.
Browser-native small molecules use Gasteiger charges instead of Sage's preferred AM1-BCC or NAGL workflows.
WebGPU uses single precision and omits periodic boundaries, particle-mesh Ewald electrostatics, barostats, and production explicit-solvent protocols.
ANI-2x accepts only a declared neutral, closed-shell element and size domain; OpenFold accepts one short chain.
Rosemary fixtures execute externally prepared protein systems but do not provide arbitrary in-browser protein parameterization.
These exclusions keep the interface from enlarging the method's claim.

Molarium separates a generated interaction layer from a versioned scientific core.
Agents may generate views, reports, or analyses; molecular state, model identity, parameters, scope checks, reference evidence, and provenance remain explicit.
Otherwise automation scales epistemic error.
Chemistry agents already join language models with literature, code, tools, and laboratory hardware~\cite{bran2024chemcrow,boiko2023coscientist}; Molarium supplies their scientific contract.

\section*{Value moves to models, intelligence, and skills}

Molarium exposes a three-layer stack.
Its MIT-licensed code, workflow state, visualizations, workers, and exports form the implementation and orchestration layer---the layer agents commoditize fastest.
OpenFF Sage, OpenMM, ANI-2x, and OpenFold form the scientific-model layer, where value rests on accuracy, domain coverage, calibration, parameterization, and prospective validation.
Molarium omits the third layer: proprietary discovery intelligence such as negative results, project history, design rationale, assay context, expert judgment, and organizational objectives.
Local Lab provides a boundary for applying that intelligence without automatically disclosing molecular structures.

\begin{figure}[htbp]
  \centering
  \begin{tikzpicture}[
    layer/.style={
      draw=MidnightBlue,
      rounded corners=2pt,
      minimum height=0.62cm,
      text width=0.78\linewidth,
      align=center,
      inner xsep=6pt,
      font=\small
    }
  ]
    \node[layer, fill=Gray!12] (implementation) at (0,0) {\textbf{Implementation and orchestration}\\interfaces, wrappers, workflow plumbing, visualization, and reports};
    \node[layer, fill=Cerulean!14] (models) at (0,0.95) {\textbf{Scientific models}\\representations, parameters, trained weights, calibration, and prospective validation};
    \node[layer, fill=ForestGreen!16] (intelligence) at (0,1.9) {\textbf{Proprietary discovery intelligence}\\experimental context, negative results, rationale, expert judgment, and organizational policy};
    \draw[->, thick, MidnightBlue]
      (7.1,0) -- (7.1,1.9)
      node[midway, right=1mm, align=center, font=\footnotesize, text=MidnightBlue] {More scarcity\\and defensibility};
    \node[font=\footnotesize, text=Gray] at (0,-0.7) {Agents most directly reduce the marginal cost and lead time of this layer.};
  \end{tikzpicture}
  \caption{Three interdependent layers of value. As implementation becomes easier to regenerate, defensible value shifts toward validated models and decision-relevant organizational intelligence.}
  \label{fig:value-layers}
\end{figure}

Implementation still matters: it makes models and organizational knowledge usable.
It is simply less scarce than the knowledge that constrains it.
Strong systems connect all three layers without blurring ownership: the interface reveals the model, the model exposes its evidence, and policy has an accountable owner.
Molarium makes the first two layers inspectable; a governed connection to the third remains the next product problem.

Proprietary intelligence is also more than a collection of assay tables and molecular structures.
It includes why a measurement was distrusted, why a chemical series stopped, which approximation failed on a target class, and how evidence is balanced against safety, novelty, timeline, and portfolio constraints.
Much of this context has historically lived in scientists' heads, project meetings, slide decks, and informal conventions.
Agents create an opportunity to apply it at the moment of decision, but only if organizations preserve provenance, disagreement, and uncertainty rather than flattening them into a single confident answer.

Agents make the third layer operational through reusable scientific skills.
A scientific skill should be more than a prompt or script.
It should define its question, domain, inputs, versions, decision rules, controls, escalation conditions, uncertainty treatment, provenance, acceptance tests, and accountable owner.
Principles for making scientific digital objects findable, accessible, interoperable, and reusable treat reuse and machine actionability as explicit goals~\cite{wilkinson2016fair}; validated skills extend that ambition from data to scientific action.

\begingroup\interlinepenalty=10000
This is where commercial value becomes more interesting rather than disappearing.
A generic feature may run a free-energy calculation.
An organization-specific skill can encode how the organization prepares systems, selects perturbations, judges convergence, handles disagreement with experiment, and decides whether a result is strong enough to change synthesis priorities.
The value lies not in the elegance of its prose but in the evidence and institutional knowledge it reliably enacts.
\par\endgroup

The human role consequently moves toward judgment.
Automation has long made operators more important at rare and consequential failure boundaries~\cite{bainbridge1983ironies}.
In molecular design, scientists still decide which question matters, which approximation is acceptable, which anomaly requires an independent control, and what consequence follows from being wrong.
Agents can make those judgments executable and can challenge them with alternatives, but accountability does not disappear into the workflow.

\section*{The product contract changes}

Fixed applications still matter for training, regulated operations, routine work, and stable tasks.
The product contract changes: generated views and workflows must retain trusted scientific components, explicit data boundaries, reproducible state, and inspectable evidence.
Molarium demonstrates this across editing, preparation, conformer comparison, minimization, dynamics, and structure prediction while keeping each method's name and scope visible.

Agents may generate presentation and orchestration, but they must not silently substitute a charge model, alter protonation, change a stopping rule, or discard a failed hypothesis.
Every result must record what ran: molecular state, model and parameter versions, external services, numerical settings, applicability limits, validation status, and human intervention.
Fluid interfaces require stable invariants.
Molarium exposes them through preparation audits, method labels, downloadable artifacts, and fail-closed scope checks.

Commercially, differentiation moves toward resources that are difficult to imitate or substitute~\cite{barney1991resources}.
It also moves toward the ability to integrate specialized knowledge distributed across an organization~\cite{grant1996knowledge}.
Companies can still charge for excellent software, but the defensible offering increasingly includes validated models, curated benchmarks, governed access to proprietary context, secure execution, scientific support, and accountable stewardship.
Composability becomes an advantage: customers should be able to create project-specific experiences without severing the evidence chain that makes the underlying methods trustworthy.

This model also creates a maintenance obligation.
Generated interfaces multiply the artifacts an organization must review, version, retire, and reproduce.
Machine-learning systems already accrue technical debt from data dependencies, configuration, and changing external conditions~\cite{sculley2015debt}.
Generated infrastructure worsens that burden when generation outpaces ownership.
Scientific skills need tests, release policy, monitoring, and named stewards.

\section*{The harder revolution}

The first agentic revolution is making computational molecular science accessible.
A scientist can increasingly request a missing interface, analysis, or workflow and receive it while the question is still alive.
Molarium demonstrates this across structure building and preparation, conformer comparison, local mechanics and neural potentials, trajectory inspection, and browser-local protein prediction.
Its strongest feature is candor: it exposes limits as clearly as capabilities.

The second revolution is harder.
It requires models that generalize to new chemistry, represent the relevant biological context, quantify uncertainty, and survive prospective experimental tests.
When anyone can generate a tool, the strategic question is no longer simply who owns the interface.
It is who owns the evidence, model, representation, skill, and judgment that make the interface worth trusting.

Software implementation is becoming cheap.
Scientific truth is not.

\section*{References}
\vspace{-0.8em}
\renewcommand{\bibsection}{}
\begin{multicols}{2}
\interlinepenalty=10000
\bibliographystyle{unsrtnat}
\bibliography{molarium_opinion_references}
\end{multicols}

\end{document}
